\chapter{Cervical Biopsy}

\section*{Overview}
A cervical biopsy is the removal of a small sample of tissue from the cervix for histologic examination. It is used to diagnose cervical abnormalities, particularly precancerous or cancerous lesions.

\section*{Alternative Names}
Colposcopy-directed biopsy, Cervical punch biopsy, Cone biopsy (conization), Loop electrosurgical excision procedure (LEEP) biopsy

\section*{Principle}
Using colposcopy, the cervix is inspected after applying acetic acid, which makes abnormal epithelium visible, and biopsy forceps sample the affected areas. The tissue is fixed and examined microscopically for dysplasia or malignancy. Larger excisional techniques such as LEEP or cone biopsy remove a conically shaped specimen that includes the transformation zone.

\section*{History}
Cervical biopsy developed with the introduction of the colposcope in the 1920s, and loop electrosurgical excision came into use in the 1980s.

\section*{Why It Is Done}
Cervical biopsy is ordered when cervical cytology (Pap smear) is abnormal, when the cervix appears abnormal on examination, or to evaluate suspicious lesions or unexplained bleeding. It confirms the presence and grade of cervical dysplasia or cancer before treatment.

\section*{Is This a Primary or Secondary Test or Procedure}
Primary diagnostic test for confirming cervical pathology identified by screening.

\section*{How It Is Performed}
With the patient in the lithotomy position, a speculum is placed and the cervix is visualized under colposcopy after application of acetic acid. Punch forceps remove small samples of abnormal tissue, or an electrosurgical loop removes a larger cone specimen. Bleeding is controlled with pressure, silver nitrate, or Monsel's solution.

\section*{Preparation}
Informed consent is obtained, and the procedure is ideally scheduled when the patient is not menstruating. The patient is advised to expect mild cramping, and analgesics may be offered.

\section*{Contraindications and Precautions}
Biopsy is avoided during active pelvic infection or acute cervicitis and in bleeding disorders that increase hemorrhagic risk. In pregnancy, biopsy of the transformation zone is generally deferred. Patients should avoid tampons, douching, and intercourse before the procedure.

\section*{Risks}
Bleeding, infection, and cramping are the main risks; cone biopsy and LEEP carry additional risks of cervical stenosis and adverse pregnancy outcomes. Minor bleeding and dark discharge may persist for several days.

\section*{Aftercare and Recovery}
The patient may resume activity but should avoid tampons, douching, and intercourse for one to two weeks, longer after LEEP or cone biopsy. Results are usually available within one to two weeks.

\section*{Results and Interpretation}
Results report the presence and grade of dysplasia (CIN 1, 2, or 3) or invasive carcinoma, including the margin status of excised specimens. Abnormal results guide treatment and follow-up, while normal results confirm no pathologic change in the sampled tissue.

\section*{Expected Results}
No dysplastic or malignant cells are present; the cervical tissue is benign.

\section*{Follow-up}
Management follows the grade of dysplasia and ranges from surveillance with repeat cytology and HPV testing to treatment with excision or ablation. Post-treatment surveillance continues for years because of the risk of recurrence.

\section*{References}
\begin{enumerate}
  \item MedlinePlus. Cervical biopsy. U.S. National Library of Medicine; 2025.
  \item Current Medical Diagnosis and Treatment. McGraw-Hill; 2025.
\end{enumerate}

\clearpage
